% =========================================================================
\section{Introducción y contexto histórico}\label{sec:introduccion}
% =========================================================================
La electrostática estudia las interacciones entre cargas eléctricas \emph{en
reposo}. Es el escalón más sencillo del electromagnetismo y, a la vez, el lugar
donde aparecen por primera vez las ideas que sostienen toda la teoría: el
\emph{campo} como intermediario de la interacción, las \emph{fuentes} que lo
generan, y la pareja de ecuaciones, una para la divergencia y otra para el
rotacional, que lo determinan por completo. Este tema parte de un único dato
experimental (la ley de Coulomb) y de un principio (la superposición), y de ellos
deduce, con las herramientas del Tema~0~\cite{TEMA0}, la ley de Gauss, el potencial
electrostático y las condiciones que cumplen los campos en una frontera. Seguimos
el programa de los apuntes de la asignatura~\cite{APUNTES-T1} y los desarrollos de
la clase, y nos apoyamos en los textos clásicos~\cite{GRIFFITHS,WANGSNESS,JACKSON,PURCELL}. Las
figuras reproducen, como diagramas en TikZ y gráficas en pgfplots, las del material
de figuras del curso~\cite{FIGURAS-T1} (preparadas con Mathematica) y las
completan con esquemas propios de la teoría y de los problemas~\cite{PROBLEMAS-T1}.

% -------------------------------------------------------------------------
\subsection{El electromagnetismo clásico y sus límites de validez}\label{sec:limites}
% -------------------------------------------------------------------------
Electromagnetismo~I y~II se enmarcan en la \emph{física clásica}. A comienzos del
siglo~XX se comprobó que la mecánica de Newton resulta «inadecuada» fuera de su
dominio natural: para objetos que se mueven a velocidades próximas a la de la luz
hay que sustituirla por la relatividad especial; para objetos muy pequeños, a escala
atómica, por la mecánica cuántica; y cuando ambas circunstancias se dan a la vez, por
la teoría cuántica de campos (\cref{fig:teorias}).

% Figura: las cuatro teorías de la física y el dominio de la electrodinámica clásica.
\begin{figure}[ht]
\centering
\begin{tikzpicture}[
    caja/.style={draw=azulclaro, fill=azulfondo, line width=0.5pt, text width=4.6cm,
                 align=center, inner sep=6pt, font=\small},
    caja0/.style={caja, draw=azulacero, line width=0.9pt},
    flecha/.style={->, line width=0.7pt, azulacero, shorten >=2pt, shorten <=2pt},
    rotulo/.style={font=\scriptsize, text=grispizarra, align=center, fill=white, inner sep=1.5pt}]
  \node[caja0] (A) at (0,0) {\textbf{Mecánica clásica}\\[1pt]{\footnotesize Newton}};
  \node[caja]  (B) at (8.8,0) {\textbf{Mecánica cuántica}\\[1pt]{\footnotesize Bohr, Heisenberg, Schrödinger}};
  \node[caja]  (C) at (0,-2.5) {\textbf{Relatividad especial}\\[1pt]{\footnotesize Einstein}};
  \node[caja]  (D) at (8.8,-2.5) {\textbf{Teoría cuántica de campos}\\[1pt]{\footnotesize Dirac, Pauli, Feynman}};
  \draw[flecha] (A) -- node[rotulo] {objetos muy pequeños\\$S\sim\hbar$} (B);
  \draw[flecha] (A) -- node[rotulo] {velocidades $v\sim c$} (C);
  \draw[flecha] (C) -- node[rotulo] {objetos muy pequeños\\$S\sim\hbar$} (D);
  \draw[flecha] (B) -- node[rotulo] {velocidades $v\sim c$} (D);
\end{tikzpicture}
\caption{Dominio de validez de las teorías de la física. El recuadro resaltado es la
física clásica, que incluye el electromagnetismo de Maxwell; este, a diferencia de la
mecánica de Newton, ya es compatible con la relatividad especial y se funde con la
cuántica en la electrodinámica cuántica.}
\label{fig:teorias}
\end{figure}


La pregunta natural es entonces si tiene sentido emplear la descripción
\emph{ondulatoria} de Maxwell. Hoy sabemos que toda radiación electromagnética
admite una descripción corpuscular (fotones) y otra ondulatoria, y que cuál de las
dos resulta más útil depende de la energía de los fotones. A frecuencias bajas, un
fotón transporta tan poca energía que cualquier detector responde a un número
enorme de ellos y solo registra su promedio: el campo se comporta como una onda
continua. Al subir la frecuencia, la naturaleza corpuscular se vuelve evidente
(por ejemplo, en los experimentos con rayos~X).

\begin{ejemplo}[Cuántos fotones hay en una onda clásica]\label{ej:fotones}
El flujo de fotones a una distancia $d$ de una fuente isótropa de potencia $P$ y
frecuencia $f$ es $\Phi_\gamma=P/(4\pi d^{2}\,hf)$, con $h=6{,}63\cdot10^{-34}\,\mathrm{J\,s}$.
\begin{itemize}
  \item \emph{Antena de FM}: $f=200\,\mathrm{MHz}$, $P=10\,\mathrm{kW}$, $d=10\,\mathrm{km}$.
        Con $hf\simeq1{,}3\cdot10^{-25}\,\mathrm{J}$ resulta
        $\Phi_\gamma\sim10^{19}$--$10^{20}\ \mathrm{fotones/(m^{2}\,s)}$.
  \item \emph{Bombilla}: $P=100\,\mathrm{W}$, luz visible ($\lambda\sim500\,\mathrm{nm}$,
        $f\sim10^{14}\,\mathrm{Hz}$, $hf\simeq4\cdot10^{-19}\,\mathrm{J}$), a $d=1\,\mathrm{m}$:
        $\Phi_\gamma\sim2\cdot10^{19}\ \mathrm{fotones/(m^{2}\,s)}$.
\end{itemize}
En ambos casos el sensor recibe una cantidad tan grande de fotones que responde a su
promedio, como lo haría ante una onda continua, y el planteamiento ondulatorio es
legítimo.
\end{ejemplo}

Dos hechos explican por qué la descripción clásica sigue siendo esencial en la física
moderna. Primero, la relatividad especial se apoyó en experimentos electromagnéticos y las
ecuaciones de Maxwell resultaron ser ya compatibles con ella. Segundo, las correcciones
cuánticas a la fuerza electromagnética son muy pequeñas incluso a distancias del orden de
$10^{-10}\,\mathrm{cm}$, cien veces menores que el tamaño de un átomo: las mismas leyes
describen la atracción y la repulsión entre partículas dentro del átomo y entre cuerpos
macroscópicos (aunque luego sea preciso recurrir a la mecánica cuántica para predecir cómo
se mueven esas partículas bajo tales fuerzas) y, a distancias aún menores, la teoría
electromagnética se funde con la cuántica en la electrodinámica cuántica. Por «clásico»
entendemos, por tanto, «no cuántico»: todas las magnitudes (posiciones, velocidades,
campos) se consideran medibles con precisión ilimitada.

% -------------------------------------------------------------------------
\subsection{La interacción electromagnética entre las cuatro fuerzas de la naturaleza}\label{sec:fuerzas}
% -------------------------------------------------------------------------
La mecánica dice cómo responde un sistema a una fuerza dada; la pregunta de fondo es de
dónde vienen las fuerzas. Todas las que conocemos se reducen a cuatro interacciones
fundamentales, que ordenamos por intensidad relativa, dando el valor~1 a la más intensa
(\cref{tab:fuerzas}).

\begin{table}[ht]
\centering
\caption{Las cuatro interacciones fundamentales. La intensidad es la relativa entre
dos protones, tomando como unidad la fuerza nuclear fuerte.}
\label{tab:fuerzas}
\small
\renewcommand{\arraystretch}{1.3}
\begin{tabular}{@{}lccl@{}}
\toprule
Interacción & Intensidad relativa & Alcance & Magnitud a la que se acopla\\
\midrule
Fuerte          & $1$          & corto ($\sim10^{-15}\,\mathrm{m}$) & carga de color\\
Electromagnética& $10^{-2}$    & largo                              & carga eléctrica\\
Débil           & $10^{-5}$    & muy corto                          & carga débil\\
Gravitatoria    & $10^{-38}$   & largo                              & masa-energía\\
\bottomrule
\end{tabular}
\end{table}

Conviene darse cuenta de la amplitud del dominio electromagnético. La fricción, la fuerza
normal que sostiene los cuerpos sobre el suelo, la que mantiene unidos los enlaces
moleculares o la que actúa en el choque de dos bolas de billar son, todas, fuerzas
electromagnéticas. La fuerza nuclear fuerte mantiene unidos los protones en el núcleo, pero
es de muy corto alcance: en un núcleo con muchos protones (como el uranio, con 92) el equilibrio es
frágil y el núcleo puede romperse. La interacción débil, responsable de ciertos procesos de
desintegración nuclear, es también de corto alcance y mucho menos intensa. La
electromagnética, asociada a la carga eléctrica, es de largo alcance y domina la escala
atómica en cuanto se supera la nuclear. La gravitatoria, asociada a la masa, es también de
largo alcance, pero extraordinariamente débil en comparación.

\begin{observacion}[Electricidad frente a gravedad]\label{obs:gravedad}
La repulsión eléctrica entre dos electrones es unas $4\cdot10^{42}$ veces mayor que su
atracción gravitatoria; el cociente $e^{2}/(4\pi\eo\,Gm_e^{2})$ no depende de la distancia, pues ambas
fuerzas decaen como $1/r^{2}$. Si la gravedad, y no la electricidad, mantuviese ligado al
electrón en el átomo de hidrógeno, el análogo del radio de Bohr,
$a_G=\hbar^{2}/(Gm_e^{2}m_p)\approx10^{29}\,\mathrm{m}$, sería mayor que el radio del
universo observable ($\sim4\cdot10^{26}\,\mathrm{m}$).
\end{observacion}

% -------------------------------------------------------------------------
\subsection{Breve historia: de la piedra de imán a la teoría de Maxwell}\label{sec:historia}
% -------------------------------------------------------------------------
Una exposición correcta del electromagnetismo pasa por conocer su historia, y el
planteamiento de este curso, que va de lo más sencillo a lo más complicado, es parecido a su
desarrollo histórico (\cref{fig:historia})~\cite{WHITTAKER,HEILBRON}.

% Figura: línea temporal (no a escala) del electromagnetismo hasta 1905.
\begin{figure}[ht]
\centering
\begin{tikzpicture}[x=1.52cm, y=1cm]
  \draw[line width=1.1pt, azulacero, -{Stealth[length=6pt]}] (-0.4,0) -- (9.5,0);
  \tikzset{ev/.style={font=\scriptsize, align=center, text width=2.55cm, text=black!85},
           anio/.style={font=\scriptsize\bfseries, text=azulprof}}
  \foreach \i/\an/\txt/\pos in {
      0/{s.\,VI a.\,C.}/{Tales de Mileto:\\el ámbar frotado}/1,
      1/{1600}/{Gilbert,\\\emph{De Magnete}}/-1,
      2/{1771}/{Cavendish:\\ley del cuadrado}/1,
      3/{1785}/{Coulomb:\\balanza de torsión}/-1,
      4/{1800}/{Volta:\\la pila}/1,
      5/{1820}/{Oersted y Ampère}/-1,
      6/{1831}/{Faraday:\\inducción y campo}/1,
      7/{1873}/{Maxwell:\\\emph{Tratado}}/-1,
      8/{1888}/{Hertz:\\ondas e.\,m.}/1,
      9/{1900--05}/{Planck,\\Einstein}/-1}{
    \fill[azulprof] (\i,0) circle (2.2pt);
    \ifnum\pos=1
      \draw[grispizarra!60, line width=0.4pt] (\i,0.1) -- (\i,0.55);
      \node[anio, anchor=south] at (\i,0.52) {\an};
      \node[ev, anchor=south] at (\i,0.95) {\txt};
    \else
      \draw[grispizarra!60, line width=0.4pt] (\i,-0.1) -- (\i,-0.55);
      \node[anio, anchor=north] at (\i,-0.52) {\an};
      \node[ev, anchor=north] at (\i,-0.95) {\txt};
    \fi}
  \node[font=\tiny, text=grispizarra, anchor=north east] at (9.45,-0.12) {no a escala};
\end{tikzpicture}
\caption{Hitos del electromagnetismo, de la Antigüedad al cambio de siglo. El eje
no está a escala: casi todo ocurre entre 1771 y 1888.}
\label{fig:historia}
\end{figure}


Las raíces de las palabras \emph{electricidad} y \emph{magnetismo} se remontan a la
Antigüedad. Tales de Mileto (siglo~VI a.\,C.) observó que el ámbar (en griego
$\ddot{\eta}\lambda\epsilon\kappa\tau\rho o\nu$, \emph{élektron}), al ser frotado, atraía
pajitas. El conocimiento del magnetismo nace de la observación de las fuerzas entre un mineral,
la «piedra de imán», y trozos de hierro; el origen de la palabra es oscuro. Lucrecio la
hace derivar de Magnesia, en Asia Menor, y Plinio refiere la leyenda del pastor Magnes, cuyo
cayado de hierro quedó prendido a una roca cercana. En 1600, William Gilbert separó con
cuidado los fenómenos eléctricos de los magnéticos en su \emph{De Magnete}.

El salto cuantitativo llega en el siglo~XVIII. En 1771 Henry Cavendish demostró, con un
experimento de esferas concéntricas (\cref{sec:coulomb-exp}), que la fuerza eléctrica
decae con el cuadrado de la distancia; sus resultados no se publicaron hasta que
Maxwell los editó en 1879~\cite{CAVENDISH-ELEC}. En 1785 Charles-Augustin de
Coulomb publicó con una balanza de torsión la ley que lleva su nombre~\cite{COULOMB}.

Electricidad y magnetismo fueron dos ciencias independientes hasta 1820, cuando Hans Christian
Oersted descubrió (casi por accidente) que una corriente eléctrica desvía una brújula: nacía el
electromagnetismo. Pocos meses después, André-Marie Ampère desarrolló una teoría matemática de
la fuerza entre elementos de corriente. Michael Faraday introdujo una idea nueva, distinta de
la acción a distancia de la mecánica newtoniana: el concepto de \emph{campo}, y en 1831
descubrió las corrientes de inducción que aparecen en conductores sometidos a un flujo magnético
variable, base de generadores, motores y transformadores. A partir de los experimentos de
Faraday, James Clerk Maxwell reunió en 1873, en su \emph{Tratado de electricidad y
magnetismo}~\cite{MAXWELL}, todo el formalismo matemático, incluidas las ecuaciones que hoy
llevan su nombre: un campo magnético variable crea un campo eléctrico sin necesidad de soporte
material, la corriente de desplazamiento y la teoría electromagnética de la luz. Su predicción de
las ondas electromagnéticas fue confirmada por Heinrich Hertz en 1888, y abrió el camino de las
telecomunicaciones (Marconi, 1896). Ni la teoría de los cuantos (Planck, 1900) ni la relatividad
especial (Einstein, 1905) han empañado el brillo de las ecuaciones de Maxwell.

% -------------------------------------------------------------------------
\subsection{De la acción a distancia al concepto de campo}\label{sec:campo-idea}
% -------------------------------------------------------------------------
La ley de Coulomb es un ejemplo de \emph{acción a distancia}: dos cargas separadas se
atraen o repelen sin que nada ocupe el espacio entre ellas. Esa idea es heredera de la gravitación
newtoniana. Sin embargo, esta interacción (como otras) se estudia con mucha más comodidad
como el encadenamiento de dos procesos: en el primero, un grupo de cargas, las \emph{fuentes},
perturban el espacio que las rodea, dotándolo de propiedades que no tenía antes; en el segundo,
otra carga sufre una fuerza neta por interactuar con ese campo previamente creado. El campo es,
pues, una propiedad de cada punto del espacio, que el teórico puede estudiar sin saber quién va
a ser la carga que lo sienta. Lo que hace a esta descripción tan fecunda es que relacionar
campos con fuentes es un problema mucho más rico que calcular fuerzas.

El plan del tema sigue una cadena lógica sencilla (\cref{fig:esquema}). La \emph{ley de
Coulomb} junto con el \emph{principio de superposición} determinan el campo eléctrico
$\vE(\vc{r})$. De la forma $1/R^{2}$ de la ley se deduce la \emph{divergencia} del campo
(ley de Gauss) y, de su carácter central, su \emph{rotacional} (nulo), de modo que $\vE$ deriva
de un potencial. Como el teorema de Helmholtz asegura que un campo queda determinado cuando se
conocen su divergencia y su rotacional (con condiciones de regularidad en el infinito), esas dos
ecuaciones, junto con las condiciones de contorno, contienen \emph{toda} la información de la
electrostática en el vacío.

% Figura: esquema lógico del tema (equivalente al esquema de los apuntes manuscritos).
\begin{figure}[ht]
\centering
\begin{tikzpicture}[
    caja/.style={draw=azulclaro, fill=azulfondo, line width=0.5pt, align=center,
                 inner sep=5pt, font=\small, minimum height=1.0cm},
    final/.style={draw=azulacero, fill=white, line width=0.8pt, align=center,
                  inner sep=5pt, font=\small, minimum height=1.0cm},
    fl/.style={->, line width=0.7pt, azulacero, shorten >=2pt, shorten <=2pt},
    nota/.style={font=\scriptsize, text=grispizarra, align=center}]
  \node[caja] (L) at (0,0) {Ley de Coulomb\\[-1pt]{\footnotesize (experimental)}};
  \node[caja] (S) at (0,-1.75) {Principio de\\superposición};
  \node[caja] (E) at (3.7,-0.87) {Campo eléctrico\\$\vE(\vc{r})$};
  \node[caja] (D) at (7.5,0.35) {$\nab\cdot\vE=\rho/\eo$\\[-1pt]{\footnotesize (ley de Gauss)}};
  \node[caja] (R) at (7.5,-2.1) {$\nab\times\vE=\vc{0}$\\[-1pt]{\footnotesize $\vE=-\nab\phi$}};
  \node[final] (H) at (11.5,-0.87) {Condiciones\\de contorno};
  \draw[fl] (L.east) -- (E.north west);
  \draw[fl] (S.east) -- (E.south west);
  \draw[fl] (E.east) -- (D.west);
  \draw[fl] (E.east) -- (R.west);
  \draw[fl] (D.east) -- (H.north west);
  \draw[fl] (R.east) -- (H.south west);
  \node[nota] at (7.5,-0.87) {teorema de\\Helmholtz};
\end{tikzpicture}
\caption{Estructura lógica del tema. De la ley de Coulomb y la superposición se deduce
el campo; su divergencia (ley de Gauss) y su rotacional (campo conservativo, potencial)
lo determinan por el teorema de Helmholtz, y, aplicadas a una frontera, dan las
condiciones de contorno.}
\label{fig:esquema}
\end{figure}


% -------------------------------------------------------------------------
\subsection{Notación y convenios del tema}\label{sec:notacion}
% -------------------------------------------------------------------------
Mantenemos la notación del Tema~0~\cite{TEMA0,NAVAS} y las expresiones del
formulario~\cite{FORMULARIO}. Los vectores llevan flecha, los versores sombrero, y con
$\vc{r}$ designamos el punto en el que se evalúa el campo (\emph{punto campo}) y con
$\vc{r}\,'$ el punto donde está la carga que lo crea (\emph{punto fuente}). Como en el
Tema~0, usamos $\sen$ para el seno.

\begin{table}[ht]
\centering
\caption{Notación específica del Tema~1.}
\label{tab:notacion1}
\small
\renewcommand{\arraystretch}{1.35}
\begin{tabular}{@{}p{4.4cm}p{9.4cm}@{}}
\toprule
Símbolo & Significado\\
\midrule
$q,\ \qp$ & Carga de prueba y carga fuente (en C).\\
$\vc{r},\ \vc{r}\,'$ & Posición del punto campo y del punto fuente.\\
$\vc{R}=\vc{r}-\vc{r}\,',\ R=\norm{\vc{R}},\ \uR=\vc{R}/R$ &
  Vector de posición relativa (de la fuente al punto campo), su módulo y su versor.\\
$\vE(\vc{r}),\ \vF$ & Campo eléctrico (V/m) y fuerza (N).\\
$\rho,\ \sigma,\ \lambda$ & Densidades volúmica (C/m$^{3}$), superficial (C/m$^{2}$) y lineal (C/m) de carga.\\
$\phi(\vc{r})$ & Potencial electrostático (V).\\
$\flujo=\oiint_{S}\vE\cdot\dS$ & Flujo eléctrico a través de una superficie cerrada.\\
$\dOm$ & Elemento de ángulo sólido (sr).\\
$\eo,\ \mu_{0},\ c$ & Permitividad y permeabilidad del vacío; velocidad de la luz.\\
$\delta(x),\ \delta^{3}(\vc{r})$ & Delta de Dirac en una y en tres dimensiones.\\
\bottomrule
\end{tabular}
\end{table}

\begin{notacion}[Convenios del tema]\label{not:convenios1}
Trabajamos en el vacío, con cargas en reposo (estática) y campos de clase $C^{1}$ salvo en
las superficies y curvas de carga, que se tratan como idealizaciones (\cref{sec:contorno}). Las
superficies cerradas se orientan con la normal exterior $\un$ y escribimos $\dS=\un\,\dsc$.
Las unidades son las del Sistema Internacional (SI).
\end{notacion}
